% !TEX program = pdflatex
\documentclass[a4paper,11pt]{article}
\usepackage[utf8]{inputenc}
\usepackage{amsmath,amssymb,amsthm}
\usepackage{mathtools}
\usepackage{booktabs}
\usepackage{longtable}
\usepackage{array}
\usepackage{hyperref}
\usepackage{geometry}
\geometry{margin=25mm}

\hypersetup{
  colorlinks=true,
  linkcolor=blue,
  urlcolor=blue,
  citecolor=blue
}

\setlength{\parindent}{1em}
\setlength{\parskip}{0.5em}

\title{Chain Complex Structure on the 4D Hypercubic Lattice: \\ Structural Correspondence between Kihara Cube-Packing \\ and Wilson Lattice Gauge Theory \\ \large (Paper 8: $\alpha$-Identity II)}
\author{Noriaki Kihara (WF System Co., Ltd.) \\ ORCID: \href{https://orcid.org/0009-0004-6753-4020}{0009-0004-6753-4020} \\ Concept DOI: \href{https://doi.org/10.5281/zenodo.19880467}{10.5281/zenodo.19880467} \\ v2 Version DOI (latest, retitled to ``Structural Correspondence''): \href{https://doi.org/10.5281/zenodo.19881119}{10.5281/zenodo.19881119}}
\date{April 2026 (v2)}

\begin{document}
\maketitle

\section*{Abstract}

This paper clarifies the mathematical positioning of the self-consistent identity $\alpha^{-1} = N(1) + V_4(1)\cdot\alpha = 137 + (\pi^2/2)\alpha$ (predicting $\alpha$ to 8.7 ppb precision) observed in the companion paper (Paper 7: $\alpha$-Identity). Paper 7 started from a purely geometric approach (4D cube-packing) but reached a precision too high to be coincidence. Reviewing the relationship to standard gauge theory, we have found that \textbf{the geometry of the 4D hypercubic lattice and Wilson lattice gauge theory (Wilson 1974) are structurally isomorphic as chain complexes at the level of chain complex, symmetry group, and discrete Stokes structure}.

In this paper, we (i) define the standard chain complex $C_\bullet = (C_0, C_1, C_2, C_3, C_4)$ on the 4D Euclidean integer lattice $\mathbb{Z}^4$, (ii) show that Schl\"afli duality (tesseract $\leftrightarrow$ 16-cell) provides the isomorphism $C_n \leftrightarrow C_{4-n}$, (iii) demonstrate that the hypercubic symmetry group $B_4$ acts equivariantly on both, and (iv) show that the discrete Stokes theorem produces identical bulk-boundary decompositions. Wilson theory is the formalism that places gauge fields on $C_1$ (links) and an action on $C_2$ (plaquettes); Kihara cube-packing is the formalism that places volumes on $C_4$ (4-cells) and measures on $C_3$ (boundary faces). They are \textbf{structurally isomorphic as chain complexes} via Schl\"afli duality.

This isomorphism opens (a) gauge-theoretic computational tools (Wilson strong-coupling expansion, renormalisation group, Monte Carlo, etc.) for application to the Kihara programme, and (b) conversely, Kihara-side number-theoretic identities (Lagrange--Jacobi four-square theorem, inclusion-exclusion) as new computational methods for Wilson theory. Paper 7's inside/outside decomposition $1 = 137\alpha + V_4(1)\alpha^2$ is naturally interpreted as a bulk-boundary normalisation condition on this chain complex.

A first-principles derivation of the numerical value of $\alpha$, the privilege of $R = 3$, and the geometric identification of the residual $c_3 \approx 1.6 \times 10^{-3}$ are left as open problems beyond this paper's scope. The contribution of this paper is to demonstrate, in standard mathematical language, that Paper 7 is not mere numerology but \textbf{a natural manifestation of the topological structure of the 4D lattice}.

\section{Introduction}

\subsection{Reach of Paper 7 and Motivation}

In the companion paper (Paper 7, Concept DOI: 10.5281/zenodo.19869266, v3), we observed the self-consistent algebraic identity
\begin{equation}
\alpha^{-1} = N(1) + V_4(1) \cdot \alpha = 137 + \frac{\pi^2}{2}\alpha \tag{1.1}
\end{equation}
constructed from $N(1) = 137$ (the count of integer-centred unit cubes fully contained in a 4D ball of radius 3) and $V_4(1) = \pi^2/2$ (the volume of the 4D unit ball). Solving the quadratic $(\pi^2/2)\alpha^2 + 137\alpha - 1 = 0$ for $\alpha$ yields $\alpha = 7.29735194\times 10^{-3}$, agreeing with the CODATA 2018 value to a relative error of $8.7 \times 10^{-8}$ (8.7 ppb).

This precision is approximately 1/3000 of the Eddington-style integer-fitting deviation ($2.6 \times 10^{-4}$), too good to be dismissed as ``mere numerical coincidence''. Paper 7 was carefully crafted to evade the Eddington trap by emphasising (i) the self-consistent equation structure, (ii) the independent 4D-geometric derivation of both 137 and $\pi^2/2$, and (iii) the QFT-perturbative analogy. However, the physical mechanism and the privilege of $R = 3$ remained open.

The motivation of this paper is the suspicion, given the strikingly high precision, that there must be a deeper mathematical structure underlying the identity. Specifically, on reviewing standard lattice gauge theory established by Wilson (1974), we recognise that the geometric setup of Paper 7 and Wilson's lattice formulation are \textbf{almost identically isomorphic from a geometric perspective}. The aim of this paper is to make this isomorphism precise in categorical language.

\subsection{Scope of the Claim}

This paper \textbf{claims} the following:

\begin{enumerate}
\item The standard chain complex on the 4D Euclidean integer lattice $\mathbb{Z}^4$ admits an equivariant action of the hypercubic symmetry group $B_4$.
\item Schl\"afli duality (the self-duality of the tesseract / 16-cell) provides the isomorphism $C_n \leftrightarrow C_{4-n}$.
\item Wilson lattice gauge theory and Kihara cube-packing are formalisms that place different structures on this chain complex, and they are \textbf{structurally isomorphic as chain complexes} via Schl\"afli duality.
\item The discrete Stokes theorem produces an identical bulk-boundary decomposition in both formalisms.
\item Paper 7's inside/outside decomposition $1 = 137\alpha + V_4(1)\alpha^2$ is naturally interpreted from this isomorphism structure.
\end{enumerate}

This paper does \textbf{not} claim:

\begin{enumerate}
\item That the value of $\alpha$ is derived from first principles via this isomorphism (we cite Paper 7 as observation only)
\item That the privilege of $R = 3$ is proved by this isomorphism
\item That U(1) gauge symmetry is uniquely emergent from this framework
\item The geometric identification of $c_3 \approx 1.6 \times 10^{-3}$
\item The specific connection to Higgs mechanism or electroweak symmetry breaking
\end{enumerate}

These are all explicitly stated as open problems in §8.

\section{Chain Complex on the 4D Hypercubic Lattice}

\subsection{Lattice Definition}

Denote the 4D Euclidean integer lattice by $\Lambda = \mathbb{Z}^4 \subset \mathbb{R}^4$. Each $x \in \Lambda$ is a \textbf{site} (vertex). Directed edges between adjacent sites are \textbf{links}, with the set of links denoted $L$. The smallest 2D faces enclosed by links are \textbf{plaquettes}, 3D cells are \textbf{3-cells}, and 4D cells are \textbf{4-cells} (tesseracts).

\subsection{Chain Complex}

Free abelian groups generated by these cells with formal integer coefficients:
\[ C_0 = \mathbb{Z}\langle\text{vertices}\rangle, \ C_1 = \mathbb{Z}\langle\text{links}\rangle, \ C_2 = \mathbb{Z}\langle\text{plaquettes}\rangle, \ C_3 = \mathbb{Z}\langle\text{3-cells}\rangle, \ C_4 = \mathbb{Z}\langle\text{4-cells}\rangle. \]

The boundary operators $\partial_n : C_n \to C_{n-1}$ are defined standardly (with orientation):
\begin{equation}
\partial_n \circ \partial_{n+1} = 0. \tag{2.1}
\end{equation}

This makes $C_\bullet$ a chain complex.

\subsection{Action of $B_4$}

The hypercubic symmetry group $B_4$ (Coxeter group of order $|B_4| = 2^4 \cdot 4! = 384$) acts on each $C_n$, equivariantly with $\partial_n$:
\begin{equation}
g \cdot \partial_n = \partial_n \cdot g, \quad \forall g \in B_4. \tag{2.2}
\end{equation}

\subsection{$f$-Vector of the Tesseract}

\begin{center}
\begin{tabular}{cccc}
\toprule
$n$ & Unit cell of $C_n$ & Boundary cells per unit & Per tesseract count \\
\midrule
0 & vertex & --- & 16 \\
1 & link & 2 vertices & 32 \\
2 & plaquette & 4 links & 24 \\
3 & 3-cell & 6 plaquettes & 8 \\
4 & 4-cell & 8 3-cells & 1 \\
\bottomrule
\end{tabular}
\end{center}

\section{Schl\"afli Duality and Poincar\'e Duality}

\subsection{Self-Duality of Tesseract and 16-cell}

Tesseract ($\{4,3,3\}$) and 16-cell ($\{3,3,4\}$) are mutually dual. The $k$-cell of the tesseract corresponds to the $(4-k-1)$-cell of the 16-cell:

\begin{center}
\begin{tabular}{ll}
\toprule
Tesseract & 16-cell \\
\midrule
16 vertices ($k=0$) & 16 cells ($k=3$) \\
32 edges ($k=1$) & 32 faces ($k=2$) \\
24 faces ($k=2$) & 24 edges ($k=1$) \\
8 cells ($k=3$) & 8 vertices ($k=0$) \\
\bottomrule
\end{tabular}
\end{center}

\subsection{Poincar\'e Duality on the Chain Complex}

For the dual lattice $\Lambda^* \cong \mathbb{Z}^4$ (half-integer-shifted), there is a discrete Poincar\'e duality:
\begin{equation}
D : C_n(\Lambda) \xrightarrow{\sim} C_{4-n}(\Lambda^*), \tag{3.1}
\end{equation}
which is $B_4$-equivariant:
\begin{equation}
D \circ g = g \circ D, \quad \forall g \in B_4. \tag{3.2}
\end{equation}

\subsection{Concrete Examples of the Dual Action}

\begin{center}
\begin{tabular}{ll}
\toprule
$\Lambda$ side & $\Lambda^*$ side \\
\midrule
Site ($C_0$) & 4-cell ($C_4$) \\
Link ($C_1$) & 3-cell ($C_3$) \\
Plaquette ($C_2$) & Plaquette ($C_2$, self-dual) \\
3-cell ($C_3$) & Link ($C_1$) \\
4-cell ($C_4$) & Site ($C_0$) \\
\bottomrule
\end{tabular}
\end{center}

The self-duality $C_2 \to C_2^*$ is the direct origin of the isomorphism between Wilson plaquette action and Kihara cube-packing.

\section{Wilson Lattice Gauge Theory in Chain-Complex Language}

\subsection{Link Variables and Gauge Transformation}

Each link $\ell = (x, y) \in L$ carries an element $U_\ell = e^{i\theta_\ell} \in U(1)$. Gauge transformation:
\begin{equation}
U_{xy} \to e^{i\alpha(x)} U_{xy} e^{-i\alpha(y)}. \tag{4.1}
\end{equation}

\subsection{Plaquette Phase and Wilson Action}

Plaquette phase:
\begin{equation}
U_p = \prod_{\ell \in \partial p} U_\ell, \tag{4.2}
\end{equation}
gauge invariant under (4.1). Wilson action:
\begin{equation}
S_{\text{Wilson}} = \frac{1}{g^2} \sum_{p \in C_2} \left[1 - \text{Re}\, U_p\right]. \tag{4.3}
\end{equation}

\subsection{Discrete Stokes Theorem}

For any closed 2D surface $\Sigma$:
\begin{equation}
\prod_{p \in \Sigma} U_p = \prod_{\ell \in \partial \Sigma} U_\ell. \tag{4.4}
\end{equation}
This is the discrete Stokes theorem, a consequence of $\partial^2 = 0$.

\section{Kihara Cube-Packing in Chain-Complex Language}

\subsection{Unit Cube Configuration}

\begin{equation}
\square(c) = \prod_{i=1}^{4} \left[c_i - \frac{1}{2}, c_i + \frac{1}{2}\right] \in C_4(\Lambda^*). \tag{5.1}
\end{equation}

A Kihara unit cube is a 4-cell of $\Lambda^*$, equivalent to a site of $\Lambda$ via $D$.

\subsection{Packing Number and Volume Deficit}

Packing condition: $\sum_{i=1}^{4} (|c_i| + 1/2)^2 \le R^2$. By direct enumeration, $N(1) = 137$ at $R=3$.

Volume deficit:
\begin{equation}
\Delta(R) = V_4(R) - N(k), \quad \Delta(R) = \frac{16\pi}{3} R^3 - 6\pi R^2 + O(R). \tag{5.2}
\end{equation}

\subsection{Bulk-Boundary Structure}

\begin{itemize}
\item \textbf{Bulk}: cubes pack exactly in the interior, contributing the strict integer $N_{\text{bulk}}$.
\item \textbf{Boundary}: cubes are partially included near $\partial B(R)$, contributing $\Delta(R)$.
\end{itemize}

This is the same structure as the discrete Stokes theorem (§4.3): bulk cancels, only boundary remains.

\section{Main Theorem: Structural Correspondence Theorem}

\subsection{Definition of $\mathcal{L}_4$}

\textbf{Definition 6.1}: $\mathcal{L}_4$ is the category of chain complexes on the 4D Euclidean integer lattice with $B_4$-equivariant boundary operators, and morphisms are $B_4$-equivariant chain maps.

\subsection{Isomorphism Theorem}

\textbf{Theorem 6.2 (Main Theorem)}: As objects of $\mathcal{L}_4$,
\[ (\mathcal{C}_W, \partial_\bullet, B_4) \cong_{\mathcal{L}_4} (\mathcal{C}_K, \partial_\bullet, B_4) \]
where $\mathcal{C}_W$ is Wilson's chain complex (gauge fields on $C_1$, action on $C_2$) and $\mathcal{C}_K$ is Kihara's (indicator on $C_4$, measure on $C_3$). The isomorphism is given by Schl\"afli duality $D$ (§3.2).

\subsection{Sketch of Proof}

(i) $\Lambda \cong \Lambda^*$ via half-integer shift, $B_4$-equivariantly.

(ii) Schl\"afli duality $D : C_n(\Lambda) \to C_{4-n}(\Lambda^*)$ is the discrete Poincar\'e duality, commuting with $\partial$.

(iii) $B_4$ acts equivariantly on both; $D$ is $B_4$-equivariant.

(iv) The discrete Stokes theorem (§4.3 and §5.3) takes the same form in both.

The detailed algebraic proof follows standard methods (Hatcher 2002). $\blacksquare$

\subsection{Corollaries}

\textbf{Corollary 6.3}: Wilson lattice gauge theory's analytical methods (strong-coupling expansion, Monte Carlo, RG) become directly applicable to the Kihara cube-packing problem via Schl\"afli duality $D$.

\textbf{Corollary 6.4}: Conversely, Kihara-side number-theoretic identities (Lagrange--Jacobi four-square theorem $r_4(N) = 8\sigma(N)$, Paper 3 §6) may provide novel closed forms for finite-volume Wilson partition functions.

\section{Mathematical Positioning of Paper 7's $\alpha$ Identity}

\subsection{Naturalness of Inside/Outside Decomposition}

Paper 7's decomposition:
\begin{equation}
1 = \underbrace{137 \alpha}_{\text{inside (bulk)}} + \underbrace{V_4(1) \alpha^2}_{\text{outside (boundary)}} \tag{7.1}
\end{equation}
is naturally interpreted from Theorem 6.2 as a bulk-boundary normalisation condition on the chain complex of the 4D lattice.

\subsection{Wilson Action Correspondence}

The formalism placing the Wilson action on Paper 7's packing problem:
\begin{equation}
S_{\text{Kihara-Wilson}} = \sum_{p \subset B(3)} [1 - \text{Re}\, U_p] \tag{7.2}
\end{equation}
is mathematically well-defined, encompassing Paper 7's geometric setup as a finite-volume version of Wilson theory.

\subsection{Repositioning of $\alpha$ Identity}

\begin{quote}
Paper 7's self-consistent identity $\alpha^{-1} = N(1) + V_4(1)\alpha$ is a structural identity naturally positioned, on the chain complex structure of the 4D Euclidean integer lattice, by the structural correspondence between Wilson lattice gauge theory and Kihara cube-packing.
\end{quote}

This is decisively different from Eddington-style numerical fitting. However, the proof that $\alpha$'s numerical value is uniquely determined from this structure is beyond this paper's scope (§8).

\section{Open Problems}

\begin{enumerate}
\item \textbf{Mechanistic derivation of $\alpha$'s value}: whether $\alpha = 1/137.036$ is uniquely determined from the isomorphism
\item \textbf{Privilege of $R = 3$}: why $N(1) = 137$ corresponds to physical $\alpha$
\item \textbf{Geometric identification of $c_3 \approx 1.6 \times 10^{-3}$}
\item \textbf{Generalisation to other gauge groups}: $SU(2), SU(3), SU(N)$, $\alpha_s$ analogues
\item \textbf{Projection to physical (3+1)D spacetime}: connection via central projection (Papers 2, 6) and Wick rotation
\item \textbf{Connection with Higgs mechanism}: implementation of the longitudinal scalar beat coupling
\item \textbf{Numerical verification}: Wilson Monte Carlo, Lagrange--Jacobi applications
\end{enumerate}

\section{Conclusion}

This paper has clarified the mathematical positioning of the self-consistent identity $\alpha^{-1} = N(1) + V_4(1)\alpha$ (predicting $\alpha$ to 8.7 ppb precision) observed in Paper 7, as the chain complex structure of the 4D Euclidean integer lattice.

The main result is that \textbf{Wilson lattice gauge theory and Kihara cube-packing are structurally isomorphic as chain complexes on the chain complex ($B_4$ symmetry, Schl\"afli duality, discrete Stokes structure) of the 4D hypercubic lattice} (Theorem 6.2). Implications:

\begin{enumerate}
\item Paper 7's inside/outside decomposition is naturally interpreted as a bulk-boundary normalisation on the chain complex
\item Wilson theory's analytical methods become applicable to the Kihara programme
\item Kihara-side number-theoretic identities may provide novel computations for Wilson theory
\item The $\alpha$ identity is not mere numerology but a natural manifestation of the topological structure of the 4D lattice
\end{enumerate}

The first-principles derivation of $\alpha$'s specific value, the privilege of $R = 3$, and the geometric identification of $c_3$ are explicitly listed as open problems in §8.

The contribution of this paper is not the mechanistic derivation of $\alpha$ but the demonstration that Paper 7 connects naturally to standard mathematical structures (algebraic topology + lattice gauge theory). This further reinforces Paper 7's evasion of the Eddington trap, and positions this research programme to be in dialogue with standard physics and mathematics communities.

\section*{References}

\begin{enumerate}
\item Wilson, K. G. ``Confinement of Quarks.'' \textit{Phys. Rev. D} \textbf{10}, 2445 (1974).
\item Kogut, J. B. ``An introduction to lattice gauge theory and spin systems.'' \textit{Rev. Mod. Phys.} \textbf{51}, 659 (1979).
\item Hatcher, A. \textit{Algebraic Topology}. Cambridge University Press (2002).
\item Coxeter, H. S. M. \textit{Regular Polytopes}. Dover Publications (1973).
\item Schl\"afli, L. \textit{Theorie der vielfachen Kontinuit\"at}. Z\"urcher und Furrer, Z\"urich (1901).
\item Bekenstein, J. D. ``Black Holes and Entropy.'' \textit{Phys. Rev. D} \textbf{7}, 2333 (1973).
\item Hawking, S. W. ``Particle Creation by Black Holes.'' \textit{Commun. Math. Phys.} \textbf{43}, 199 (1975).
\item Tangherlini, F. R. ``Schwarzschild Field in $n$ Dimensions and the Dimensionality of Space Problem.'' \textit{Nuovo Cimento} \textbf{27}, 636 (1963).
\item Eddington, A. S. ``The Charge of an Electron.'' \textit{Proc. Roy. Soc. London} \textbf{A122}, 358 (1929).
\item Mohr, P. J., Newell, D. B., Taylor, B. N. ``CODATA Recommended Values of the Fundamental Physical Constants: 2018.'' \textit{Rev. Mod. Phys.} \textbf{93}, 025010 (2021).
\end{enumerate}

\section*{Appendix A: $f$-Vector and Schl\"afli Duality}

Tesseract ($\{4,3,3\}$): $(f_0, f_1, f_2, f_3) = (16, 32, 24, 8)$.
16-cell ($\{3,3,4\}$): $(f_0, f_1, f_2, f_3) = (8, 24, 32, 16)$.
Schl\"afli duality: $f_k(\text{Tesseract}) = f_{3-k}(\text{16-cell})$ (Coxeter 1973).

\section*{Appendix B: Numerical Reproduction of Paper 7's $\alpha$ Identity (reference)}

Identical to Paper 7 §3.2 (v3, 8.7 ppb precision):

\begin{verbatim}
V_4(1) = π²/2 ≈ 4.9348022
N(1) = 137
Equation: 4.9348 α² + 137 α - 1 = 0
Positive root: α ≈ 7.29735194 × 10⁻³
α⁻¹ ≈ 137.0360110
CODATA 2018: α⁻¹ = 137.0359991
Relative error: 8.7 × 10⁻⁸ (8.7 ppb)
\end{verbatim}

For details, see Paper 7 (Concept DOI: 10.5281/zenodo.19869266).

\end{document}
